% Exact scalar dependency used in PDE9--11 and IPM.20--21.
\section{The positive Gamma grid product in the original period determinant}
\label{sec:PGCC}

The full original PDE determinant proof reduces its constant to
\[
 \prod_{r=1}^{d-1}\Gamma(r/d)
       =\frac{(2\pi)^{(d-1)/2}}{\sqrt d},\qquad d=n+1\geq2.
 \tag{PGCC1}
\]
Here is a complete proof of this particular scalar identity from the
positive defining Gamma integral, retaining its positive square root.
For a real $0<a<1$, Tonelli's theorem applies to the nonnegative
integrand defining the product of two Gamma integrals. Set
$x=st/(1+t)$ and $y=s/(1+t)$, with $s,t>0$. This is a bijective
change of variables on the positive quadrant with Jacobian
$s/(1+t)^2$, and direct multiplication gives
\[
 \begin{split}
 \Gamma(a)\Gamma(1-a)
 &=\int_0^\infty\int_0^\infty
        e^{-x-y}x^{a-1}y^{-a}\,dx\,dy\\
 &=\int_0^\infty e^{-s}\,ds
       \int_0^\infty\frac{t^{a-1}}{1+t}\,dt
  =\int_0^\infty\frac{t^{a-1}}{1+t}\,dt.
 \end{split}\tag{PGCC2}
\]
The last integral converges at zero because $a>0$ and at infinity
because $a<1$.

To evaluate it with its exact sign, use
$z^{a-1}=\exp((a-1)\operatorname{Log}z)$ with $0<\arg z<2\pi$.
The positively oriented keyhole contour about the positive real axis
has an upper segment from $\varepsilon$ to $R$, an outer circle
counterclockwise, a lower segment from $R$ to $\varepsilon$, and an
inner circle clockwise. Take $0<\varepsilon<1<R$ and let the banks
approach the cut. The only pole of $z^{a-1}/(1+z)$ inside is $z=-1$,
with residue $e^{\pi i(a-1)}$. The two straight segments sum to
\[
 (1-e^{2\pi ia})\int_\varepsilon^R
                       \frac{t^{a-1}}{1+t}\,dt.
 \tag{PGCC3}
\]
The outer circular contribution has absolute value at most
$2\pi R^a/(R-1)$, which tends to zero as $R\to\infty$.
The inner contribution has absolute value at most
$2\pi\varepsilon^a/(1-\varepsilon)$, which tends to zero as
$\varepsilon\to0$. The residue theorem and these bounds give
\[
 \begin{split}
 (1-e^{2\pi ia})\int_0^\infty\frac{t^{a-1}}{1+t}\,dt
 &=2\pi i e^{\pi i(a-1)},\\
 \Gamma(a)\Gamma(1-a)&=\frac{\pi}{\sin(\pi a)}.
 \end{split}\tag{PGCC4}
\]
In the second line the quotient sign follows explicitly from
$1-e^{2\pi ia}=-2i e^{\pi ia}\sin(\pi a)$ and
$e^{\pi i(a-1)}=-e^{\pi ia}$.

Now put $\zeta=e^{2\pi i/d}$. Factoring $X^d-1$ over its distinct
roots and evaluating the quotient by $X-1$ at $X=1$ proves
\[
 d=\prod_{r=1}^{d-1}(1-\zeta^r),\qquad
 d=\prod_{r=1}^{d-1}|1-\zeta^r|
   =2^{d-1}\prod_{r=1}^{d-1}\sin(\pi r/d).
 \tag{PGCC5}
\]
The sines in this product are all positive. Apply PGCC4 at
$a=r/d$ and multiply over $1\leq r<d$. The involution
$r\mapsto d-r$ permutes the same factors, so
\[
 \left(\prod_{r=1}^{d-1}\Gamma(r/d)\right)^2
 =\prod_{r=1}^{d-1}\frac{\pi}{\sin(\pi r/d)}
 =\frac{(2\pi)^{d-1}}d.
 \tag{PGCC6}
\]
Each Gamma integral is strictly positive. Taking the positive square
root proves PGCC1, including the case $d=2$.

This evaluates the precise scalar used in PDE9--11 without replacing
the original Fourier determinant or its phase. In fact the polynomial
period isomorphism already follows before this evaluation: PDE9 gives
nonzero positive Gamma factors and fixed nonzero powers of $u$, while
PDE10 gives the nonzero Fourier Vandermonde determinant. PDE6--8 then
propagate that nonzero determinant through the full coefficient family.
PGCC1 additionally supplies its displayed value $(2\pi)^{n/2}$,
with all original powers and phases calculated in PDE9--11 unchanged.
